% polyglossia's Japanese and Chinese: \XeTeXlinebreaklocale "ja" and "zh" with
% \XeTeXlinebreakskip, ICU's line breaks (in Latin Modern: the CJK glyphs
% are missing, as in the oracle; the breaks do not depend on them).
\documentclass{article}
\usepackage{polyglossia}
\setmainlanguage{english}
\setotherlanguage{japanese}
\setotherlanguage{chinese}
\setmainfont{Latin Modern Roman}
\newfontfamily\japanesefont{Latin Modern Roman}
\newfontfamily\chinesefont{Latin Modern Roman}
\showboxdepth=100 \showboxbreadth=100000
\begin{document}
English text first.
\textjapanese{「こんにちは」、世界。日本語の文章です。句読点、括弧（かっこ）や「かぎ括弧」も。ちょっと待ってください……ねえ、ゃゅょっァィゥェォ。東京は２０２６年１０月６日（火）です。ラーメン・カレー・ソバ〜！？}
Back to English.

\begin{japanese}
吾輩は猫である。名前はまだ無い。どこで生れたかとんと見当がつかぬ。何でも薄暗いじめじめした所でニャーニャー泣いていた事だけは記憶している。
\end{japanese}

\textchinese{这是一个测试：中文、英文（English）和数字123混排。《三国演义》第一回——宴桃园豪杰三结义。}

\setbox0\vbox{\hsize=80pt \textjapanese{吾輩は猫である。名前はまだ無い。}}\showbox0
\end{document}
